% Created 2026-05-25 Mon 06:38
% Intended LaTeX compiler: pdflatex
\documentclass[11pt]{article}
\usepackage[utf8]{inputenc}
\usepackage[T1]{fontenc}
\usepackage{graphicx}
\usepackage{longtable}
\usepackage{wrapfig}
\usepackage{rotating}
\usepackage[normalem]{ulem}
\usepackage{amsmath}
\usepackage{amssymb}
\usepackage{capt-of}
\usepackage{hyperref}
\author{Santiago Arango-Piñeros}
\date{2026-02-12}
\title{Math 411 - Spring 2026\\\medskip
\large Introduction to Abstract Algebra 1}
\hypersetup{
 pdfauthor={Santiago Arango-Piñeros},
 pdftitle={Math 411 - Spring 2026},
 pdfkeywords={},
 pdfsubject={},
 pdfcreator={Emacs 29.4 (Org mode 9.6.15)}, 
 pdflang={English}}
\begin{document}

\maketitle
\tableofcontents


\section*{Coordinates}
\label{sec:org0545789}
\subsection*{Instructor}
\label{sec:org7ea8ca7}
\href{https://sarangop1728.github.io/}{Santiago Arango-Piñeros}, Ph.D. (He/Him).


My office is in \textbf{LGRT 1238}, and you can contact me via email at:
\begin{itemize}
\item \texttt{sarangopiner@umass.edu}.
\end{itemize}

\subsection*{Teaching Assistants}
\label{sec:orgcf889b9}
\begin{itemize}
\item Natalie Welling (\texttt{nwelling@umass.edu})
\item Sam Silver (\texttt{samuelsilver@umass.edu})
\end{itemize}

\subsection*{Lectures}
\label{sec:org45fc277}
\begin{center}
\begin{tabular}{l|l|l|l}
\hline
Section & Days & Time & Place\\[0pt]
\hline
02-LEC (80294) & Mo, We & 4:00 PM - 5:15 PM & LGRT 171\\[0pt]
03-LEC (85426) & Mo, We & 5:30 PM - 6:45 PM & LGRT 171\\[0pt]
\hline
\end{tabular}
\end{center}

\subsection*{Office hours}
\label{sec:org41a2627}

If you are planning on coming to my office hours, I'd appreciate an email
letting me know that you are planning to attend.

\begin{center}
\begin{tabular}{l|l|l|l}
\hline
Days & Time & Place & By\\[0pt]
\hline
Mo, We & 3:00 PM - 3:50 PM & LGRT 1238 & Santiago\\[0pt]
Fri & 2:00 PM - 3:50 PM & LGRT 1522 & Natalie and Sam\\[0pt]
\hline
\end{tabular}
\end{center}


If the times above are not convenient, or if you would like to discuss a
private matter, send me an email.

\section*{Course description}
\label{sec:org916bda7}
The focus of the course will be on studying \emph{groups}. These are algebraic
structures that capture the notion of symmetry. Groups are ubiquitous in all
areas of mathematics and the world around us (see \hyperref[sec:org7bb381f]{these links}). If you commit
to this class, you will master the essential concepts of group theory outlined
in the \hyperref[sec:org2102cd8]{learning objectives} by the end of the course.
\subsection*{Prerequisites}
\label{sec:org278d2a4}
\textbf{MATH 235} and either \textbf{CMPSCI 250} or \textbf{MATH 300}. In other words, we will need
some important concepts from linear algebra, and there will be an emphasis on
proofs and the development of precise mathematical reasoning and writing.

\subsection*{Textbook}
\label{sec:orgba38f80}
\begin{itemize}
\item \textbf{\emph{Algebra: Abstract and Concrete}} by Frederick M. Goodman. The book is
freely available for download at the \href{https://homepage.divms.uiowa.edu/\~goodman/algebrabook.dir/algebrabook.html}{author's website}.
\item As a complement to the textbook, we will use some of \href{https://kconrad.math.uconn.edu/blurbs/}{Keith Conrad's blurbs} on
group theory.
\item Every textbook has strengths and weaknesses. You might find it useful to
complement your reading with the following resources:
\begin{itemize}
\item Aluffi, \href{https://bookstore.ams.org/gsm-104}{\emph{Algebra: Chapter 0}}.
\item Dummit and Foote, \href{https://books.google.com.co/books/about/Abstract\_Algebra.html?id=KJDBQgAACAAJ\&redir\_esc=y}{\emph{Abstract Algebra}}.
\item Macauley, \href{https://www.math.clemson.edu/\~macaule/visualalgebra.html}{\emph{Visual Algebra}}.
\item Terras, \href{https://books.google.com.co/books/about/Abstract\_Algebra\_with\_Applications.html?id=znWIDwAAQBAJ\&redir\_esc=y}{\emph{Abstract Algebra with Applications}}.
\item Rotman, \href{https://books.google.com.co/books/about/A\_First\_Course\_in\_Abstract\_Algebra.html?id=ctEZAQAAIAAJ\&redir\_esc=y}{\emph{A First Course in Abstract Algebra}}.
\end{itemize}
\end{itemize}

\subsection*{Learning objectives}
\label{sec:org2102cd8}

\begin{enumerate}
\item To learn the fundamental examples of groups. In this class, these are:
\begin{itemize}
\item The finitely generated abelian groups. In particular, the cyclic groups
\(\mathbb{Z}\) and \(\mathbb{Z}/n\mathbb{Z}\).
\item The dihedral groups \(D_{n}.\)
\item The symmetric group \(S_{n}.\)
\item The alternating group \(A_{n}.\)
\item The matrix groups \(\mathrm{SO}_2(\mathbb{R}), \mathrm{SO}_3(\mathbb{R}),
     \mathrm{GL}_2(R)\), and \(\mathrm{SL}_2(R)\) for \(R = \mathbb{Z}/n\mathbb{Z}, \mathbb{R}\).
\item The affine groups \(\mathrm{Aff}_2(R)\) for \(R =
     \mathbb{Z}/n\mathbb{Z}, \mathbb{R}\).
\end{itemize}

\item To learn the axiomatic definition of a group and how to use it to prove
basic properties.
\item To learn the concepts of subgroups, cosets, quotients, and how to combine
these to derive Lagrange’s theorem.
\item To learn the concepts of homomorphisms and isomorphisms and Noether's
isomorphism theorems.
\item To understand the structure theorem of finitely generated abelian groups in
terms of the Smith normal form of a matrix with integer coefficients.
\item To learn what it means for a group to act on a set as well as the
natural actions of each of the fundamental examples.
\item To understand how to use the Sylow theorems.
\end{enumerate}

\subsection*{Homework}
\label{sec:org06b116b}
Homework assignments are a key component of this class. Even though they won't
be graded, the midterm exams will be based on the homework problems. \uline{It is in
your best interest to invest considerable time thinking about the homework
problems.}

Collaboration with other students is highly encouraged! Come to office hours to
work on problems with other students.

\begin{itemize}
\item \texttt{PSET0}: Back to basics (\href{https://www.dropbox.com/scl/fi/9cq8i56k0c4l8zxx5vrtt/pset0.pdf?rlkey=g8g80z674b1ywms2kkcddybg4\&dl=0}{\texttt{pdf}})
\item \texttt{PSET1}: Some examples of groups (\href{https://www.dropbox.com/scl/fi/t1cjodyterp1oi2o4xy1r/pset1.pdf?rlkey=p1spc9zhe18088o06lj5nu0e4\&dl=0}{\texttt{pdf}})(\href{https://www.dropbox.com/scl/fi/c7bftc22bv2utfx8neqsc/pset1-sol.pdf?rlkey=1kfofy7f0njfstujkcjyckuc6\&dl=0}{\texttt{solution}})
\item \texttt{PSET2}: Basic properties of groups (\href{https://www.dropbox.com/scl/fi/vmtjxwrqpgkuptwz8lywx/pset2.pdf?rlkey=nkuqr7ux74kakw2ffh57i7qax\&dl=0}{\texttt{pdf}})(\href{https://www.dropbox.com/scl/fi/zv651hjibk52ferk4sd7v/pset2-sol.pdf?rlkey=vgvuik1w9apj73wp7cbh26lsu\&dl=0}{\texttt{solution}})
\item \texttt{PSET3}: Cosets and Lagrange's theorem (\href{https://www.dropbox.com/scl/fi/qq8t6s1cnw7cqcjg6hgdk/pset3.pdf?rlkey=opzttzka6v2it3tg3w5rgflnu\&dl=0}{\texttt{pdf}})(\href{https://www.dropbox.com/scl/fi/v56wvste1kscud3fols0d/pset3-solution.pdf?rlkey=9wj2ndb083zenxw0iv5vuwiw9\&dl=0}{\texttt{solution}})
\item \texttt{PSET4}: The isomorphism theorems (\href{https://www.dropbox.com/scl/fi/f3163uprv0fxm4thy72hk/pset4.pdf?rlkey=9b1peazeg6l131ohybogzuerc\&dl=0}{\texttt{pdf}})(\href{https://www.dropbox.com/scl/fi/hee9mm14r98g89cv7lpaw/pset4-solution.pdf?rlkey=5pehee7boftvtxehidq8ahcag\&dl=0}{\texttt{solution}})
\item \texttt{PSET5}: Direct and semidirect products (\href{https://www.dropbox.com/scl/fi/wr5i7ci600fwxcuzdpjtb/pset5.pdf?rlkey=doky4q7bjl9gtmo03iu3774x1\&dl=0}{\texttt{pdf}})
\item \texttt{PSET6}: Finitely generated abelian groups (\href{https://www.dropbox.com/scl/fi/2pyhr73z9d4z0ev9wcoj0/pset6.pdf?rlkey=hy2qxhe7d8ldtd8dplopfip2a\&dl=0}{\texttt{pdf}})
\item \texttt{PSET7}: Group actions and the Sylow theorems (\href{https://www.dropbox.com/scl/fi/9j61gn6zdeohtkmzvr5sf/pset7.pdf?rlkey=tdc0c2wo3bhsw2u6d9llwr9jf\&dl=0}{\texttt{pdf}})
\end{itemize}

\subsection*{Midterm exams (\texttt{600 points})}
\label{sec:org3a54746}
There will be three midterm exams, each one worth \texttt{200 points}. Each exam will
have 5 questions, each one worth \texttt{40 points}. \textbf{Question 1} will ask you to
define a concept or state an important theorem. \textbf{Question 2} will ask you to
prove a result (of reasonable difficulty) from the corresponding lectures.
\textbf{Question 3} will consist of five \emph{True} or \emph{False} questions, with no
justification needed. \textbf{Question 4} will ask you to perform a computation, and
\textbf{Question 5} will be a new problem closely related to the corresponding
\texttt{PSETs}. The midterms will be held on Wednesday evenings.

\begin{itemize}
\item \texttt{EXAM1}: Emphasis on lectures 1–7 and \texttt{PSETs} 0, 1, 2. \textit{<2026-02-25 Wed 19:00>}.
Room: \texttt{LGRC A301}. (\href{https://www.dropbox.com/scl/fi/5idb8pu63uolnypv7mgof/exam1v1solution.pdf?rlkey=rpr12umksg62hd2d426a94zuo\&dl=0}{\texttt{solution}})
\item \texttt{EXAM2}: Emphasis on lectures 6–10 and \texttt{PSETs} 0, 3, 4. \textit{<2026-03-25 Wed 19:00>}.
Room: \texttt{LGRC A301}. (\href{https://www.dropbox.com/scl/fi/685wheuhivq2iiyaerdb1/solution-e2v2.pdf?rlkey=denrm1nqf6g2lj4wkf26i3xkd\&dl=0}{\texttt{solution}})
\item \texttt{EXAM3}: Emphasis on lectures 11–15 and \texttt{PSETs} 0, 5, 6.
\textit{<2026-04-15 Wed 19:00>}. Room: \texttt{GSMN 51}. (\href{https://www.dropbox.com/scl/fi/karhbg31bksh91n3iq76o/exam3v1-solution.pdf?rlkey=ng6wg3pnabgg40g20nzu22dfw\&dl=0}{\texttt{solution}})
\end{itemize}

See the \hyperref[sec:org3afed50]{``Mistakes were made'' policy} below.

\subsection*{Final exam (\texttt{200 points})}
\label{sec:org3be0c72}
This exam evaluates the contents of all lectures and all \texttt{PSETs}, with a
special focus on lectures 17–20 and \texttt{PSETs} 0, 7. It will consist of four
questions, each one worth \texttt{50 points}.
\begin{itemize}
\item \texttt{FINAL}: \textit{<2026-05-12 Tue 08:00> } AM. Room: \texttt{LGRT 123}.
\end{itemize}

\subsection*{Worksheets (\texttt{200 points})}
\label{sec:orgb6a35bd}
Sometimes, there will be a graded in-class worksheet that evaluates your
understanding of the topics covered during the lecture. I will explain the
solutions at the end of the class, and grade the worksheets only for
completion. If you don't turn in your worksheet, you will get a score of zero.
There are no make-up worksheets, but the 3 lowest scores will be dropped.

\section*{Grades}
\label{sec:orge9f43d9}
The perfect final grade is \texttt{1000 points}. To compute your grade, add the scores
obtained in your midterm exams, final exam, and worksheets, and consult the
following tables.

\begin{center}
\begin{tabular}{l|lllll}
\textbf{Grade} & A & A\(-\) & B\(+\) & B & B\(-\)\\[0pt]
\hline
\texttt{points} & \([930,1000]\) & \([900,930)\) & \([870,900)\) & \([830,870)\) & \([800,830)\)\\[0pt]
\end{tabular}
\end{center}

\begin{center}
\begin{tabular}{l|llllll}
\textbf{Grade} & C\(+\) & C & C\(-\) & D\(+\) & D & F\\[0pt]
\hline
\texttt{points} & \([770,800)\) & \([730,770)\) & \([700,730)\) & \([670,700)\) & \([600,670)\) & \([0,600)\)\\[0pt]
\end{tabular}
\end{center}

\subsection*{Mistakes were made policy}
\label{sec:org3afed50}
With the intention of rewarding behaviors that promote learning, students may
submit corrections to \texttt{EXAM1}, \texttt{EXAM2}, and \texttt{EXAM3} to \uline{earn back up to 80\% of
the points lost}.

\subsubsection*{Prerequisite: submit a summary}
\label{sec:orgc19dafe}
To qualify to submit corrections, a student must submit a signed handwritten
summary of the topics assessed on the corresponding exam. The summary must be
uploaded to \textbf{Gradescope} before the scheduled exam.
\begin{itemize}
\item Students taking the exam at a different time, for whatever reason, must still
submit their summaries before the original scheduled exam.
\end{itemize}

\subsubsection*{Submitting corrections}
\label{sec:org31c1b1b}
Corrections are due exactly two weeks after the exam date and must be submitted
through \textbf{Gradescope} before midnight.

The correction for each problem must include the following.

\begin{enumerate}
\item A complete and correct explanation of your mistake.
\item A complete and correct explanation of a solution.
\item A discussion of what you can do to prevent making this type of mistake again
in the future.
\end{enumerate}

\subsubsection*{Grading corrections}
\label{sec:orgc056562}
If the corrections comply with all the requirements listed above, they will be
accepted and the score for the corresponding question on the exam will be
adjusted. If the grader of the correction deems your original answer a
\emph{reasonable attempt}, then you will get back 80\% of the points. Otherwise, you
will get back 50\% of the points.

If the corrections fail to comply with \textbf{any} of the requirements listed above
(e.g., no summary was presented, or one of items (1), (2), or (3) was not
answered to the grader’s satisfaction), the corrections will be rejected and
the score for the corresponding question on the exam will \textbf{not} be adjusted.

\section*{Topics and schedule}
\label{sec:orgd7528bb}
It is the student's responsibility to read the material before the lecture.
During the lectures, we will focus on reviewing the key concepts, answering
questions, and working on examples.

\begin{center}
\begin{tabular}{l|l|r}
\hline
Date & Lecture & Reading\\[0pt]
\hline
\textit{<2026-02-02 Mon>} & \href{https://www.dropbox.com/scl/fi/p1wxilymw0vyhqvnk3z0a/Lecture-1.pdf?rlkey=wz6pvevnvakryws6uz4t8jsgx\&st=oihmxpsu\&dl=0}{1. What is symmetry?} & 1.1 - 1.7, \href{https://kconrad.math.uconn.edu/blurbs/grouptheory/whygroups.pdf}{Blurb}\\[0pt]
\textit{<2026-02-04 Wed>} & \href{https://www.dropbox.com/scl/fi/ye5ewjg4w53js9i9nnlqp/Lecture-2.pdf?rlkey=en0vp8k0mxcdldak8rpebqe5w\&st=q1vmzgj1\&dl=0}{2. Examples of groups} & 1.1 - 1.7\\[0pt]
\textit{<2026-02-09 Mon>} & \href{https://www.dropbox.com/scl/fi/5fn0wkb6ejtsk79mqsjpc/Lecture-3.pdf?rlkey=2v43ltzmasna2gk3lyqmmoulk\&dl=0}{3. Abstract groups: first results} & 1.10, 2.1\\[0pt]
\textit{<2026-02-11 Wed>} & \href{https://www.dropbox.com/scl/fi/17ji5counqyvroky0zxh7/Lecture-4.pdf?rlkey=ihh68st77av0gy3322151n05e\&dl=0}{4. Subgroups and cyclic groups} & 2.2\\[0pt]
\textit{<2026-02-16 Mon>} & \textbf{No class} (Presidents' Day, moved to Thursday!) & \\[0pt]
\textit{<2026-02-18 Wed>} & \href{https://www.dropbox.com/scl/fi/utckf9w72zu37yaec6f8c/Lecture-5.pdf?rlkey=fe58spmzyrto2y4xmrhr1lbh3\&dl=0}{5. Dihedral groups} & 2.3\\[0pt]
\textit{<2026-02-19 Thu>} & \href{https://www.dropbox.com/scl/fi/c4bwzgfb771pt55lv5r7k/Lecture-6.pdf?rlkey=66e9htb7zs3mo0avx0w6bcezy\&dl=0}{6. Homomorphisms and isomorphisms} & 2.4\\[0pt]
\textit{<2026-02-23 Mon>} & \href{https://www.dropbox.com/scl/fi/4z4br5snua92kfmkyyqha/Lecture-7.pdf?rlkey=red8npkhxxxf51qcykp1rudpy\&dl=0}{7. The sign of a permutation} & \href{https://kconrad.math.uconn.edu/blurbs/grouptheory/sign.pdf}{Blurb}\\[0pt]
\textit{<2026-02-25 Wed>} & \textbf{Exam 1} at 7:00 PM in \texttt{LGRC A301}. (No class) & \\[0pt]
\textit{<2026-03-02 Mon>} & \href{https://www.dropbox.com/scl/fi/c49t3qqjya7v7pc29tazj/Lecture-8.pdf?rlkey=5qmji57c31u5npkxthnpqr0ya\&dl=0}{8. Cosets} & 2.5\\[0pt]
\textit{<2026-03-04 Wed>} & \href{https://www.dropbox.com/scl/fi/oxt143yka9abuh63wmpfd/Lecture-9.pdf?rlkey=iw757fm7u5t6e1jr5u2xlqy6h\&dl=0}{9. Lagrange's theorem} & 2.5\\[0pt]
\textit{<2026-03-09 Mon>} & \href{https://www.dropbox.com/scl/fi/5umx9xqjvzhxw28wvecpl/Lecture-10.pdf?rlkey=7o83fhinujrhqwrodnnmg0mco\&dl=0}{10. Noether's first isomorphism theorem} & 2.7\\[0pt]
\textit{<2026-03-11 Wed>} & 11. Noether's other isomorphism theorems & 2.7\\[0pt]
\textit{<2026-03-16 Mon>} & \textbf{No class} (spring recess) & \\[0pt]
\textit{<2026-03-18 Wed>} & \textbf{No class} (spring recess) & \\[0pt]
\textit{<2026-03-23 Mon>} & \href{https://www.dropbox.com/scl/fi/yogahfmerlrevbc65y30r/Lecture-11.pdf?rlkey=w71u05arfw3kzyqiehc1ufx7d\&dl=0}{12. Direct products} & 3.1\\[0pt]
\textit{<2026-03-25 Wed>} & \textbf{Exam 2} at 7:00 PM in \texttt{LGRC A301}. (No class) & \\[0pt]
\textit{<2026-03-30 Mon>} & \href{https://www.dropbox.com/scl/fi/jhwp7l0ykhttzsk74jsiq/Lecture-13.pdf?rlkey=89ua250me6pwwy68oqcl830t8\&dl=0}{13. Semidirect products} & 3.2\\[0pt]
\textit{<2026-04-01 Wed>} & \href{https://www.dropbox.com/scl/fi/qi7ztpadgkzfiei306ep2/Lecture-14.pdf?rlkey=7gsry779qqv2hq7xdxn48ebsm\&dl=0}{14. Linear algebra over the integers} & 3.5\\[0pt]
\textit{<2026-04-06 Mon>} & \href{https://www.dropbox.com/scl/fi/fsajexd86jw6dg0bhtoyk/Lecture-14.pdf?rlkey=ihw92clet9umsnutu9reqh7qm\&dl=0}{15. Finitely generated abelian groups} & 3.6\\[0pt]
\textit{<2026-04-08 Wed>} & \href{https://www.dropbox.com/scl/fi/2zfxo0q6uf7v2pwq5xnqs/Lecture-16.pdf?rlkey=9c7xhkeiv50d3fnq67wz5tqtm\&dl=0}{16. Elementary divisors decomposition} & 3.6\\[0pt]
\textit{<2026-04-13 Mon>} & (Problems day) & \\[0pt]
\textit{<2026-04-15 Wed>} & \textbf{Exam 3} at 7:00 PM in \texttt{GSMN 51}. (No class) & \\[0pt]
\textit{<2026-04-20 Mon>} & \textbf{No class} (Patriot's Day) & \\[0pt]
\textit{<2026-04-22 Wed>} & \href{https://www.dropbox.com/scl/fi/5qe7hwh5bejcstqb8wihe/Lecture-18.pdf?rlkey=vggwpqmhxmvyfyq5adlhdykk0\&dl=0}{17. Group actions} & 5.1\\[0pt]
\textit{<2026-04-24 Fri>} & \href{https://www.dropbox.com/scl/fi/5e0ct4ph7dcczk7dbdblx/Lecture-19.pdf?rlkey=ai8rn33avfbgkxzsaq0udtfq7\&dl=0}{18. Counting orbits} & 5.2\\[0pt]
\textit{<2026-04-27 Mon>} & \href{https://www.dropbox.com/scl/fi/b1lgkyzgax41va73m8dcy/Lecture-20.pdf?rlkey=2wntwn600iv1rs23altpxdm88\&dl=0}{19. Symmetries of groups} & 5.3\\[0pt]
\textit{<2026-04-29 Wed>} & \href{https://www.dropbox.com/scl/fi/t8plgi2v2iy9i1xb38pco/Lecture-21.pdf?rlkey=swne53r5vpy94gxow1zqolqy8\&dl=0}{20. Group actions and group structure} & 5.4\\[0pt]
\textit{<2026-05-04 Mon>} & 21. The Sylow theorems (proofs) & \href{https://kconrad.math.uconn.edu/blurbs/grouptheory/sylowpf.pdf}{Blurb}\\[0pt]
\textit{<2026-05-06 Wed>} & 22. (Problems day) & \\[0pt]
\hline
\end{tabular}
\end{center}

\section*{Group theory links}
\label{sec:org7bb381f}
Here are some interesting Wikipedia links to various topics (in math, culture,
and science) related to group theory. I'd be happy to chat about these during
office hours.
\begin{itemize}
\item \href{https://en.wikipedia.org/wiki/Rubik\%27s\_Cube\_group}{The Rubik’s Cube Group}
\item \href{https://en.wikipedia.org/wiki/15\_puzzle}{The 15-Puzzle}
\item \href{https://en.wikipedia.org/wiki/Tower\_of\_Hanoi}{The Towers of Hanoi}
\item \href{https://en.wikipedia.org/wiki/Molecular\_symmetry}{Molecular Symmetry}
\item \href{https://en.wikipedia.org/wiki/Islamic\_geometric\_patterns}{Islamic Geometric Patterns}
\item \href{https://en.wikipedia.org/wiki/Diffie\%E2\%80\%93Hellman\_key\_exchange}{The Diffie--Hellman Key Exchange}
\item \href{https://en.wikipedia.org/wiki/Binary\_Golay\_code}{The binary Golay code}
\item \href{https://en.wikipedia.org/wiki/Free\_group}{Free Groups}
\item \href{https://en.wikipedia.org/wiki/Galois\_theory}{Galois theory}
\item \href{https://en.wikipedia.org/wiki/Fundamental\_group}{The Fundamental Group}
\item \href{https://en.wikipedia.org/wiki/Braid\_group}{Braids and Braid Groups}
\item \href{https://en.wikipedia.org/wiki/Triangle\_group}{Triangle Groups}
\item \href{https://en.wikipedia.org/wiki/Platonic\_solid}{Platonic Solids}
\item \href{https://en.wikipedia.org/wiki/Elliptic\_curve\#The\_group\_law}{Elliptic Curves}
\item \href{https://en.wikipedia.org/wiki/Monster\_group}{The Monster Group}
\end{itemize}

\section*{Philosophy}
\label{sec:org3796ae3}
\subsection*{Adopt a growth mindset}
\label{sec:org103280e}
Your effort and attitude determine your abilities. Embrace challenges and
failure as an opportunity to grow. Find inspiration in other people's success.

\subsection*{Learning is the student's responsibility}
\label{sec:org64bfe04}
Paraphrasing Galileo:
\begin{quote}
``You cannot teach a person \textbf{anything}; you can only help
them find it within themselves.''
\end{quote}
We are all here to \uline{understand}. My job as a more experienced learner is to
assist you on your journey. But you are responsible for investing the time and
effort necessary to learn.

\subsection*{Doing hard things}
\label{sec:orga33472a}
This is hard work, and it will be frustrating at times. In my opinion, the
reward is well worth the investment, as is often the case with challenging
endeavors. In the words of JFK:
\begin{quote}
``We choose to go to the Moon in this decade and do the other things, not
because they are easy, but because they are hard; because that goal will serve
to organize and measure the best of our energies and skills, because that
challenge is one that we are willing to accept, one we are unwilling to
postpone, and one we intend to win, and the others, too.''
\end{quote}

\subsection*{Everyone belongs in this classroom}
\label{sec:org2c28e79}
We will subscribe to \href{https://www.ams.org/publications/journals/notices/201610/rnoti-p1164.pdf}{Federico's axioms}.

\begin{itemize}
\item \textbf{Axiom 1.} Mathematical potential is equally present in different groups,
irrespective of geographic, demographic, and economic boundaries.

\item \textbf{Axiom 2.} Everyone can have joyful, meaningful, and empowering mathematical
experiences.

\item \textbf{Axiom 3.} Mathematics is a powerful, malleable tool that can be shaped and
used differently by various communities to serve their needs.

\item \textbf{Axiom 4.} Every student deserves to be treated with dignity and respect.
\end{itemize}

\section*{Administrative details}
\label{sec:org8c84445}
\begin{itemize}
\item Add/drop only through SPIRE.
\item I do not keep a waiting list, and the mathematics department staff will not
handle these matters.
\item Final exams are kept by the mathematics department. Copies are available upon
request.
\end{itemize}

\subsection*{Drops, withdrawals, and incompletes}
\label{sec:org85495ec}
\begin{itemize}
\item Last day to add or drop with no record: \textit{<2026-02-04 Wed>}.
\item Last day to drop with W: \textit{<2026-04-02 Thu>}.
\item See the \href{https://www.umass.edu/registrar/academic-calendar}{academic calendar} for other important dates.
\item Incomplete grades are warranted only if a student is passing the course at the
time of the request and if the course requirements can be completed by the
end of the following semester. Read more \href{https://www.umass.edu/natural-sciences/advising/petitions-and-forms/incomplete-grade-form}{here}.
\end{itemize}

\subsection*{Make-up exam policy}
\label{sec:org1e50294}
You must take the regular exam unless you qualify for an official make-up exam
approved by me, following the official make-up request procedure. Make sure you
read and understand the make-up exam procedure.

\begin{itemize}
\item \textbf{Final exam conflict:} If you need a make-up exam due to a final exam
schedule conflict, you must submit documentation from the Registrar's Office
or other supporting documents at least two weeks before the scheduled exams.
No exceptions will be made. No later than one week before the exam, you must
submit a written request to me that includes: your name and UMass Amherst
Student ID number, your section number, and the reason for requesting the
make-up exam. You can request make-up exams through your SPIRE account: in
SPIRE, go to Student Home > Final Exam Conflict.

\item \textbf{Religious observance:} If you must miss an exam due to religious observance,
you must contact me within two weeks of the beginning of the semester.

\item \textbf{Medical reasons:} If you will be absent from an exam due to medical reasons,
you must notify me at least one week in advance of the exam. If you have a
medical emergency, you must notify me as soon as possible. In either case,
you may need to provide documentation. You do not need to disclose personal
details of your condition, but you must provide enough information to allow
the absence to be excused.

\item \textbf{Other circumstances:} It is impossible to anticipate all possible
situations. In the event of an exceptional circumstance not covered above,
you must contact me and explain the problem. You must be prepared to provide
a written statement if necessary. I will evaluate the reasons you provide and
make a decision.

\item Note that there is \textbf{no re-taking of exams} in this course. If you are sick
and take the exam anyway, you cannot re-take the exam later for a better
grade. Regardless of the situation, if you do not feel you can take the exam
on the scheduled date, you must inform me as soon as possible.

\item Make-up exams will \textbf{not} be given to accommodate travel plans.

\item I will ensure that taking a make-up exam does not represent any technical
advantage. In particular, the questions will be completely different from
those on the main exam.
\end{itemize}

\subsection*{Class attendance policy}
\label{sec:org3da35fd}
By UMass policy, students are expected to attend all regularly scheduled
classes at the University for which they are registered. When planning for the
tests and homework, I will assume that you have been following my lectures.
That being said, I will not enforce or grade attendance. If you are not
able to attend one of the lectures, make sure you read the notes for that
day and talk to other students to check if you missed anything important.

\subsection*{Class etiquette}
\label{sec:org256f103}
\begin{itemize}
\item I expect you to be present and refrain from using your phone.
\item Arrive on time. If you arrive late, try to minimize your disruption.
\item Laptops and tablets are allowed during lectures, provided that you do not
disrupt your fellow classmates or the lecture.
\end{itemize}

\subsection*{Academic dishonesty}
\label{sec:org0df2147}

Academic dishonesty includes but is not limited to: 
\begin{itemize}
\item \textbf{Cheating:} intentional use, and/or attempted use of trickery, artifice,
deception, breach of confidence, fraud, and/or misrepresentation of one's
academic work.
\item \textbf{Fabrication:} intentional and unauthorized falsification and/or invention of
any information or citation in any academic exercise.
\item \textbf{Plagiarism:} knowingly representing the words or ideas of another as one's
own work in any academic exercise. This includes submitting without citation,
in whole or in part, prewritten term papers of another or the research of
another, including but not limited to commercial vendors who sell or
distribute such materials.
\item \textbf{Facilitating dishonesty:} knowingly helping or attempting to help another
commit an act of academic dishonesty, including substituting for another in
an examination, or allowing others to represent as their own one's papers,
reports, or academic works.
\end{itemize}

Formal definitions of academic dishonesty, examples of various forms of
dishonesty, and the procedures which faculty must follow to penalize dishonesty
are detailed on the \href{https://www.umass.edu/studentsuccess/academic-integrity}{Academic Honesty website}. Appeals must be filed within ten
days of notification by the Academic Honesty Office that a formal charge has
been filed by an instructor who suspects dishonesty. Contact the Academic
Honesty Office for more information on the process. The \href{https://www.umass.edu/ombuds/}{Ombuds Office} is also
available to support individuals engaging with the Academic Honesty process.
The \href{https://www.umass.edu/provost/}{Provost’s Office} is where appeals are processed and filed.

\section*{Required statements}
\label{sec:org02bb353}
\subsection*{Academic honesty statement}
\label{sec:org25f0635}
Since the integrity of the academic enterprise of any institution of higher
education requires honesty in scholarship and research, academic honesty is
required of all students at the University of Massachusetts Amherst. Academic
dishonesty is prohibited in all programs of the University. Academic dishonesty
includes but is not limited to: cheating, fabrication, plagiarism, and
facilitating dishonesty. Appropriate sanctions may be imposed on any student
who has committed an act of academic dishonesty. Instructors should take
reasonable steps to address academic misconduct. Any person who has reason to
believe that a student has committed academic dishonesty should bring such
information to the attention of the appropriate course instructor as soon as
possible. Instances of academic dishonesty not related to a specific course
should be brought to the attention of the appropriate department head or chair.
Since students are expected to be familiar with this policy and the commonly
accepted standards of academic integrity, ignorance of such standards is not
normally sufficient evidence of lack of intent
(\url{http://www.umass.edu/dean\_students/codeofconduct/acadhonesty/}).

\subsection*{Academic integrity statement}
\label{sec:orgd144d09}
UMass Amherst is strongly committed to academic integrity, which is defined as
completing all academic work without cheating, lying, stealing, or receiving
unauthorized assistance from any other person, or using any source of
information not appropriately authorized or attributed. As a community, we hold
each other accountable and support each other’s knowledge and understanding of
academic integrity. Academic dishonesty is prohibited in all programs of the
University and includes but is not limited to: cheating, fabrication,
plagiarism, lying, and facilitating dishonesty, via analogue and digital means.
Sanctions may be imposed on any student who has committed or participated in an
academic integrity infraction. Any person who has reason to believe that a
student has committed an academic integrity infraction should bring such
information to the attention of the appropriate course instructor as soon as
possible. All students at the University of Massachusetts Amherst have read and
acknowledged the Commitment to Academic Integrity and are knowingly responsible
for completing all work with integrity and in accordance with the policy:
(\url{https://www.umass.edu/senate/book/academic-regulations-academic-integrity-policy}).

\subsection*{Accommodation statement}
\label{sec:org90ab4eb}
The University of Massachusetts Amherst is committed to providing an equal
educational opportunity for all students. If you have a documented physical,
psychological, or learning disability on file with Disability Services (DS),
you may be eligible for reasonable academic accommodations to help you succeed
in this course. If you have a documented disability that requires an
accommodation, please notify me within the first two weeks of the semester so
that we may make appropriate arrangements. For further information, please
visit Disability Services (\url{https://www.umass.edu/disability/}).

\subsection*{Title IX statement}
\label{sec:orge7923ba}
In accordance with Title IX of the Education Amendments of 1972 that prohibits
gender-based discrimination in educational settings that receive federal funds,
the University of Massachusetts Amherst is committed to providing a safe
learning environment for all students, free from all forms of discrimination,
including sexual assault, sexual harassment, domestic violence, dating
violence, stalking, and retaliation. This includes interactions in person or
online through digital platforms and social media. Title IX also protects
against discrimination on the basis of pregnancy, childbirth, false pregnancy,
miscarriage, abortion, or related conditions, including recovery. There are
resources on campus to support you. A summary of the available Title IX
resources (confidential and non-confidential) can be found at the following
link: \url{https://www.umass.edu/titleix/resources}. You do not need to make a formal
report to access them. If you need immediate support, you are not alone. Free
and confidential support is available 24 hours a day / 7 days a week / 365 days
a year at the SASA Hotline 413-545-0800.
\end{document}
